Deep-learning models for stock-return prediction have mostly been developed on the markets
they were designed for, and their predictive power is seldom checked when the model is
moved to a different market. This thesis studies that cross-market transfer in one
concrete case. The model is Stockformer, a network of roughly one million parameters that
combines a wavelet decomposition, attention and a graph of stock relations, calibrated on
Chinese A-share stocks. The question is twofold: whether the model keeps its predictive
power when retrained on the S\&P 500 and, if it does not, where the net return observed in
that market comes from.

The first part compares Stockformer with simpler models (regularised regressions and
boosted trees) under a common evaluation framework, on a common data panel and at a daily
horizon. The main metric is the rank information coefficient (IC), which measures the
correlation between the predicted and the realised ordering of returns. Possible
information leakage is also audited; the only incident found favours the complex model and
is reported separately. In this implementation and on the universe analysed (477 S\&P 500
constituents with complete price coverage), Stockformer shows no robust evidence of
transferability. Its out-of-sample IC is $-0{,}003$, indistinguishable from zero, while a
five-coefficient Lasso regression reaches $+0{,}0238$. All the simple models come out
ahead of Stockformer, although by small margins.

The second part studies the origin of the net return of a weekly market-neutral strategy
on the same universe. A stage-by-stage attribution shows that the signal produces a weak
gross return and that portfolio construction with explicit transaction-cost control is the
stage that preserves it in net terms. A signal search with criteria fixed in advance, and
validated on a reserved sample, points to residual momentum as the most promising
candidate. Combined with that construction, the net Sharpe ratio rises from $0{,}38$ to
$0{,}72$ over a 311-week window and keeps its sign on the reserved sample. The improvement
is presented as the most promising result of the thesis, pending confirmation with more
data.

The conclusions are conditional on the study design: a universe restricted to current
index constituents with complete coverage, public data without historical reconstruction,
a single run of the model and a loss function adapted from the original. Under those
conditions, the contribution is twofold: empirical evidence of absence of transfer in this
case, and a reproducible evaluation and attribution procedure that places the origin of
the net return in a simple signal and in portfolio construction, not in the complexity of
the predictor.